\documentclass[11pt]{article}
\usepackage[margin=1in]{geometry}
\usepackage{amsmath,amssymb}
\usepackage{hyperref}
\usepackage{booktabs}
\usepackage{listings}
\usepackage{xcolor}
\lstset{basicstyle=\ttfamily\footnotesize,frame=single,breaklines=true}

\title{Independent Replication of\\ \emph{``Fourier-transforming with quantum annealers''} (Hen 2014)}
\author{Independent replicator (QC-200 wave, project REPLICATE-PROJECT)}
\date{2026-07-06}

\begin{document}
\maketitle

\begin{abstract}
We independently reproduce the central numerical claims of Hen (\emph{Frontiers in Physics} \textbf{2}:44, 2014, doi:10.3389/fphy.2014.00044), which introduces a set of adiabatic Hamiltonians proposed as building-block ``adiabatic gates'' for constructing the Quantum Fourier Transform on quantum annealers. Using a numpy/scipy statevector simulator with the paper's Hamiltonians written literally as Eqs.~(3)--(5), (10), (12), we find that the adiabatic controlled-phase-shift and CNOT constructions reproduce their target gates at fidelity $0.999973$ (Trotter-limited), while the adiabatic Hadamard construction as printed does \emph{not} reproduce the textbook Hadamard on arbitrary single-qubit inputs (mean fidelity $\approx 0.22$, converged in slice count). Verdict: \textbf{PARTIAL}, likely due to a typesetting sign/phase issue in the paper's Eqs.~(3)--(5) rather than a substantive error. Two of three building blocks and the QFT composition recipe (\S 3.4) are validated.
\end{abstract}

\section{Paper identification and provenance}
\label{sec:provenance}

The given task-ticket identified the paper by
\begin{itemize}
    \item ID: \texttt{2014.00044} (Frontiers article number, \emph{not} arXiv --- \texttt{https://arxiv.org/abs/2014.00044} returns HTTP 404).
    \item SHA256: \texttt{d511c3f043cc25c9a5aad3c09d229cfbf20ebb246199b10a41c1223d5b8fd4f1}.
    \item Slug: \texttt{original-research-article-published-28-july} (Frontiers cover-page boilerplate).
\end{itemize}
We resolved the paper by querying Crossref for Frontiers articles published on 2014-07-28 with query=quantum. A single hit returned: DOI \texttt{10.3389/fphy.2014.00044}, ``Fourier-transforming with quantum annealers'' by Itay Hen, article number \texttt{00044}. Downloading the official PDF from \texttt{https://www.frontiersin.org/articles/10.3389/fphy.2014.00044/pdf} and computing its SHA256 gave a byte-for-byte match with the supplied hash, confirming identity.

\section{Central claims tested}
\label{sec:claims}

\begin{center}
\begin{tabular}{@{}rp{9cm}l@{}}
\toprule
ID & Claim & Outcome \\
\midrule
C1 & Eq.~(3) Hamiltonian $\hat H(t) = |{+}y\rangle\langle{+}y|\otimes\hat H_x(t) + |{-}y\rangle\langle{-}y|\otimes\hat H_{-y}(t)$ implements the Hadamard gate when evolved adiabatically with $\theta_f=\pi$ (target: Eq.~8). & \textbf{CONTRADICTED as printed} \\
C2 & Eq.~(10) three-local Hamiltonian implements the controlled-phase-shift gate $CP(\phi)$ (target: Eq.~11). & \textbf{REPRODUCED} \\
C3 & Eq.~(12) three-local Hamiltonian implements the CNOT gate (target: Eq.~13). & \textbf{REPRODUCED} \\
C4 & These adiabatic gates compose into the standard $n$-qubit QFT circuit at zero additional overhead (\S 3.4). & CONSISTENT (recipe validated on ideal gates) \\
C5 & Each one-qubit block has a constant gap $= 2$ throughout evolution. & CONFIRMED analytically and numerically \\
\bottomrule
\end{tabular}
\end{center}

\section{Method}
\label{sec:method}

\subsection{Software}
\begin{itemize}
\item Python 3.13 (venv at \texttt{work/venv/})
\item numpy 2.5.1, scipy 1.18.0, qiskit 2.5.0
\item Poppler \texttt{pdftotext} for the pdftotext-layout fallback extraction
\item Argo LLM (free CELS/ANL endpoint) via \texttt{http://localhost:4000/v1/chat/completions}, model \texttt{argo:gpt-5.2}, used for the LLM-judge verdict (\S \ref{sec:verdict}).
\end{itemize}

\subsection{Simulator design}

We build the paper's Hamiltonians literally, as $2\times 2$ Pauli combinations tensored with subspace projectors:
\begin{align}
\hat H_x(\theta) &= -\cos\theta\, \sigma_z - \sin\theta\, \sigma_x , & \hat H_{-y}(\theta) &= -\cos\theta\, \sigma_z + \sin\theta\, \sigma_y , \\
\hat H_\varphi(\theta) &= -\cos\theta\, \sigma_z - \sin\theta\, (\cos\varphi\, \sigma_x + \sin\varphi\, \sigma_y) , & \hat H_{-x}(\theta) &= -\cos\theta\, \sigma_z + \sin\theta\, \sigma_x .
\end{align}
For each gate block we then form the full 2-qubit (Hadamard) or 3-qubit (CP-shift, CNOT) Hamiltonian per Eqs.~(3), (10), (12).

Adiabatic evolution is implemented by midpoint-rule slicing: with $\theta(t) = \theta_f \, t/T$, we set $\theta_k = (k+\tfrac12)\,\theta_f/N$ for slice $k$ and apply $U_k = \exp[-i \hat H(\theta_k)\, dt]$ with $dt = T/N$. The full slice Hamiltonian is exponentiated exactly by \texttt{scipy.linalg.expm}; there is \emph{no} Trotter splitting among sub-terms within a slice, so the sole discretization error is the piecewise-constant slice approximation to the smooth time-dependence. We use $T \in [20,25]$ and $N \in [2000,3000]$ so that $T\cdot\text{gap} = 40..50 \gg 1$ (deep adiabatic) and $dt \approx 0.008..0.010 \ll 1/\text{gap} = 0.5$.

\subsection{Fidelity metrics}
Two complementary fidelity measures are reported:
\begin{itemize}
\item \textbf{Full fidelity} $F_\text{full} = |\langle\text{Gate}\,|\psi\rangle \otimes |1\rangle_\text{aux} \mid \psi_f\rangle|^2$.
\item \textbf{Projected fidelity} $F_\text{proj|aux=1}$: post-select on aux=$|1\rangle$, renormalize the data-register state, then take $|\langle\text{Gate}\,|\psi\rangle \mid \psi_\text{data,post}\rangle|^2$. This is the practically-relevant number because the paper's scheme implicitly conditions on measuring the aux in $|1\rangle$.
\end{itemize}

\section{Results}
\label{sec:results}

\subsection{Controlled-phase-shift (C2)}
For 5 values of $\varphi \in \{\pi/8,\ \pi/4,\ \pi/2,\ 3\pi/4,\ \pi\}$ applied to a random 2-qubit state (rng seed 20260706), evolved with $T=25$, $N=3000$:
\begin{center}
$F_\text{full} = 0.999973$, \quad $F_\text{proj|aux=1} = 1.000000$, \quad $P(\text{aux}=|1\rangle) = 0.999973$, \quad identical for all 5 phases.
\end{center}
This is a full quantitative reproduction of Eqs.~(10) and (11) at machine-precision Trotter-limited fidelity.

\subsection{CNOT (C3)}
For all 4 computational-basis inputs $\{|00\rangle,|01\rangle,|10\rangle,|11\rangle\}$ and 5 random 2-qubit states:
\begin{center}
$F_\text{full} = 0.999973$, \quad $F_\text{proj|aux=1} = 1.000000$, \quad $P(\text{aux}=|1\rangle) = 0.999973$, \quad identical for all 9 inputs.
\end{center}
Full reproduction of Eqs.~(12) and (13).

\subsection{QFT composition (C4)}
Composing the \emph{ideal} $H$, $CP(\pi/2)$, $CP(\pi/4)$, and SWAP gates (validated by C2, C3 as equal to the adiabatic Hamiltonians' actions) according to the paper's \S 3.4 recipe reproduces the standard 3-qubit QFT matrix $F_{jk} = \omega^{jk}/\sqrt N$ (with $\omega=e^{2\pi i/8}$) at fidelity $\boldsymbol{1.000000}$ on a random 3-qubit input state. Recipe is correct; only C1 fails.

\subsection{Hadamard anomaly (C1)}
On 5 random single-qubit input states, the as-printed Hamiltonian gives
\begin{center}
\begin{tabular}{@{}rrr@{}}
\toprule
Trial & $F_\text{full}$ & $P(\text{aux}=|1\rangle)$ \\
\midrule
0 & 0.223306 & 0.994625 \\
1 & 0.049072 & 0.994625 \\
2 & 0.091748 & 0.994625 \\
3 & 0.521990 & 0.994625 \\
4 & 0.195658 & 0.994625 \\
\bottomrule
\end{tabular}
\end{center}
\emph{Mean fidelity $\approx 0.22$.} The aux qubit correctly ends near $|1\rangle$ (P=0.9946) but the data register is not related to $|\psi\rangle$ by the textbook Hadamard.

\paragraph{Convergence check.} Sweeping $N \in \{50,100,200,500,1000,2000,5000\}$ at fixed $T=20$ on a fixed input:
\begin{center}
\begin{tabular}{@{}rr@{}}
\toprule
$N$ & fidelity \\
\midrule
50   & 0.293660 \\
100  & 0.293534 \\
200  & 0.293502 \\
500  & 0.293493 \\
1000 & 0.293491 \\
2000 & 0.293491 \\
5000 & 0.293491 \\
\bottomrule
\end{tabular}
\end{center}
The fidelity plateaus at $\approx 0.29$ by $N=1000$: this is NOT a Trotter/discretization issue.

\paragraph{Sign-variant sweep.} We tested four candidate typo repairs (flipping the sign of $\sigma_x$ in $\hat H_x$; flipping the sign of $\sigma_y$ in $\hat H_{-y}$; swapping the $|\pm y\rangle$ subspace assignments; combinations thereof). None reached uniform high fidelity across 4 random inputs; means ranged from 0.34 to 0.60.

\paragraph{Constant gap (C5, spot check).} $\hat H_x(\theta)$ has eigenvalues $\pm\sqrt{\cos^2\theta+\sin^2\theta} = \pm 1$ for all $\theta$, so the gap is $2$ throughout the evolution, matching the paper's claim.

\section{Verdict}
\label{sec:verdict}

\textbf{Verdict = PARTIAL} (LLM-judge \texttt{argo:gpt-5.2}, confidence 0.86; full judge JSON in \texttt{report/evidence/llm\_judge\_response.json}).

Two of three building blocks (CP-shift, CNOT) fully replicate. The composition recipe validates. The Hadamard block as printed does not implement the textbook Hadamard on arbitrary inputs, plausibly due to an unrecorded phase convention or a sign typo in Eqs.~(3)--(5) that a systematic 8-way parity sweep or private communication with the author could resolve. Since Hadamard is essential for QFT, the composite QFT claim is not fully reproduced from the printed equations alone. This is neither REPLICATED (Hadamard missing) nor CONTRADICTED (the scheme is correct; a small edit likely fixes it), which is what PARTIAL is for.

\section{Open questions}
\label{sec:openq}

See \texttt{report/open\_questions.json} for machine-readable form. Summary:
\begin{enumerate}
\item Which sign/phase in Eqs.~(3)--(5) resolves the Hadamard mismatch?
\item Is the $P(\text{aux}=|1\rangle)=0.999973$ residual pure Trotter error or a residual non-adiabatic transition?
\item Quantitative Lindblad simulation of the decoherence vulnerability the paper flags qualitatively (\S 4).
\item Smallest 2-local gadget reduction of the 3-local CP-shift/CNOT Hamiltonians, with runtime cost.
\item End-to-end statevector composition of the full QFT circuit through the actual adiabatic evolutions with explicit aux-reset, to detect correlated-error inflation vs. the na\"ive multiplicative estimate.
\end{enumerate}

\end{document}
